\section{Cache Results}
\label{sec:results}

Table~\ref{tab:totals} lists total time and peak working set.
Total time includes table construction on the cached runs.
At \(n = 8\) that construction is the whole run: 1.179\,s against 0.009\,s for the original field.
At \(n = 64\) the cached total drops below the original, 1.563\,s against 1.962\,s.
At \(n = 256\) the totals are 25.3\,s and 120.1\,s.

\begin{table}[!t]
\caption{Total time and peak working set. A dash means that size was not run for that build.}
\label{tab:totals}
\centering
\scriptsize
\begin{tabular}{@{}rrrrr@{}}
\toprule
\(n\) & Cached (s) & Original (s) & Cached (MB) & Original (MB) \\
\midrule
8   & 1.179  & 0.009  & 158.6 & 4.9 \\
32  & ---    & 0.248  & ---   & 4.6 \\
64  & 1.563  & 1.962  & 159.4 & 5.4 \\
128 & ---    & 15.062 & ---   & 7.9 \\
256 & 25.291 & 120.142 & 172.3 & 18.2 \\
512 & 194.403 & ---   & 213.9 & --- \\
\bottomrule
\end{tabular}
\end{table}

The arithmetic excluding table construction is the matrix product plus the three determinants.
That portion is 4.94 times faster at \(n = 64\) (0.394\,s cached, 1.948\,s original) and 4.98 times faster at \(n = 256\) (24.11\,s cached, 120.03\,s original).
Table~\ref{tab:parts} splits the \(256 \times 256\) run.
The product is about half of the arithmetic, and the three determinants are nearly equal, which matches one elimination each on \(A\), \(B\), and \(C = AB\).
Cached time grew by a factor of 7.7 from \(n = 256\) to \(n = 512\), close to the factor of 8 expected from \(n^3\).

\begin{table}[!t]
\caption{Section times for \(n = 256\), in seconds.}
\label{tab:parts}
\centering
\scriptsize
\begin{tabular}{@{}lrr@{}}
\toprule
Section & Cached & Original \\
\midrule
Generate \(A\) and \(B\) & 0.014 & 0.107 \\
Product \(C = AB\) & 12.298 & 56.971 \\
\(\det(A)\) & 3.937 & 21.027 \\
\(\det(B)\) & 3.939 & 21.036 \\
\(\det(C)\) & 3.937 & 20.993 \\
\bottomrule
\end{tabular}
\end{table}

\begin{figure}[!t]
\centering
\begin{tikzpicture}
\begin{axis}[
  ieee lines,
  ymode=log,
  ymin=0.005,
  ylabel={Total time (s)},
  xlabel={Matrix dimension \(n\)},
  symbolic x coords={8,64,256},
]
\addplot coordinates {(8,1.179) (64,1.563) (256,25.291)};
\addplot coordinates {(8,0.008522) (64,1.962) (256,120.142)};
\legend{Cached, Original}
\end{axis}
\end{tikzpicture}
\caption{Total runtime, including table construction. The vertical axis is logarithmic so the \(8 \times 8\) original run remains visible.}
\label{fig:runtime}
\end{figure}

Peak memory tells a different story from the published laptop figures.
The cached process sits near 159\,MB from \(n = 8\) through \(n = 64\), which is the addition and multiplication maps, and rises to 214\,MB at \(n = 512\) as the matrices grow.
The original process is about 5\,MB until the matrices matter, and 18\,MB at \(n = 256\).
On this implementation the cache adds memory.
It does not shrink the matrices.
Figure~\ref{fig:memory} shows the sizes that both builds ran.

\begin{figure}[!t]
\centering
\begin{tikzpicture}
\begin{axis}[
  ieee bars,
  ymin=0,
  ylabel={Peak working set (MB)},
  xlabel={Matrix dimension \(n\)},
  symbolic x coords={8,64,256},
]
\addplot coordinates {(8,158.57) (64,159.43) (256,172.26)};
\addplot coordinates {(8,4.85) (64,5.35) (256,18.18)};
\legend{Cached, Original}
\end{axis}
\end{tikzpicture}
\caption{Peak working set on the same runs as Fig.~\ref{fig:runtime}.}
\label{fig:memory}
\end{figure}

Extending the cached \(n^3\) trend from the 194\,s run at \(n = 512\) puts a \(4000 \times 4000\) cached run near 26 hours on this CPU.
The same extrapolation of the original \(n = 256\) run puts the uncached \(4000 \times 4000\) case near 127 hours.
The assignment notes of about 20 minutes cached and 6 hours 40 minutes original were measured on an Apple M3 Max \cite{russell2026assignment}.
Those notes are a different machine and a different table layout.
They are not a prediction for the map-based tables timed here.
